%% ============================================================
%% RIFT Phase 2 — Paper I: Unlocking the Action
%% Theoretical Framework for Dark Energy and Dark Matter
%% from First Principles
%% ============================================================
\documentclass[12pt,a4paper]{article}

\usepackage{amsmath,amssymb,amsfonts}
\usepackage{geometry}
\usepackage{hyperref}
\usepackage{booktabs}
\usepackage{graphicx}
\usepackage{xcolor}

\geometry{margin=2.5cm}

\hypersetup{
  colorlinks=true,
  linkcolor=blue!60!black,
  citecolor=blue!60!black,
  urlcolor=blue!60!black
}

\newcommand{\Psi}{\Psi}
\newcommand{\pl}{M_{\mathrm{Pl}}}
\newcommand{\Lz}{\Lambda_0}
\newcommand{\half}{\tfrac{1}{2}}
\newcommand{\rift}{\textsc{rift}}
\newcommand{\lcdm}{$\Lambda$CDM}

\title{\textbf{RIFT Phase 2 — Paper I}\\[6pt]
       Unlocking the Action:\\
       Dark Energy and Dark Matter from\\
       First Principles in Recursive\\
       Intelligence-Field Theory}

\author{C.~R.~B.~Wilson}

\date{April 2026}

% ============================================================
\begin{document}
\maketitle

\begin{abstract}
Recursive Intelligence-Field Theory (RIFT) Phase~1 established a
\emph{locked action} --- a single-parameter scalar--tensor gravity with
coupling $\Lambda({\Psi}) = \Lz\Psi^2$ and source $R$ --- that passes
all 23 cosmological tests (Sims~82--111) and is observationally
indistinguishable from \lcdm{} at current precision.  Two structural
gaps remain: (i) dark energy is an unexplained cosmological constant
inserted by hand, and (ii) dark matter has no origin in the field
equations.  Phase~1 simulations explicitly closed every kinematic
alternative within the locked action (Sims~99--100, 103, 111),
establishing that both gaps \emph{require unlocking the action}.

This paper derives the unique Phase~2 action from the symmetry
principles of RIFT: diffeomorphism invariance, a $\mathbb{Z}_2$
symmetry $\Psi\to-\Psi$, second-order field equations (Horndeski
class), and the \emph{recursion principle} that $\Psi$ is both
sourced by and sources spacetime curvature.  These constraints,
together with UV finiteness established in Phase~1 (Sims~102--106),
uniquely determine the extension:

\begin{enumerate}
\item A quartic Jordan-frame potential $V_J = \lambda\Psi^4$, which
  generates --- via the conformal transformation to the Einstein frame
  --- a plateau potential that acts as an effective, \emph{derived}
  cosmological constant.  The bare $\Lambda$ is removed; dark energy
  emerges from the geometry of field space.
\item A conformal matter coupling $\beta\Psi^2\mathcal{L}_m$, which
  gives $\Psi$ an environment-dependent effective mass.  In galaxies,
  where local curvature $R_{\rm gal}$ drives a tachyonic instability,
  $\Psi$ condenses into a halo whose profile is determined by the
  \emph{recursion fixed-point condition} --- a self-consistency
  requirement that the condensate sources the curvature that sustains
  it.  This condition uniquely determines $\beta$ in terms of $\Lz$
  and observed rotation curve amplitudes.
\end{enumerate}

The Phase~2 action introduces two new parameters: $\lambda$ (DE
scale) and $\beta$ (DM coupling).  Both are determined by
observational boundary conditions, not free choices.  The programme
of Phase~2 simulations (SIM112 onward) will test whether this
action simultaneously reproduces dark energy, dark matter, and all
Phase~1 cosmological tests.
\end{abstract}

\tableofcontents
\newpage

% ============================================================
\section{Introduction}
\label{sec:intro}

\subsection{Phase 1 summary}

RIFT Phase~1 (Papers~I--V, master compilation; Sims~82--111)
established a complete programme of tests for the locked action

\begin{equation}
S_{\rm lock} = \int d^4x\,\sqrt{-g}\left[
  \frac{\pl^2 + 2\Lz\Psi^2}{2}\,R
  -\half(\nabla\Psi)^2
  -\half m_0^2\Psi^2
\right] + S_{\rm matter}
\label{eq:locked}
\end{equation}

with $\Lz = 0.003$, $m_0 = 0.001\,H_0$.  Key results:

\begin{itemize}
  \item All BAO, CMB, RSD, and Solar System tests: PASS
    ($\Delta\chi^2 = 0$ vs \lcdm{} in every dataset).
  \item UV structure: finite through two-loop graviton sector; $\Lz$
    asymptotically free (Sims~102--106).
  \item Inflation: locked action fails ($n_s = 0.932$, $>5\sigma$;
    Sim~110).  Starobinsky attractor requires UV coupling $\xi\sim10^4$.
  \item Dark energy: $w_0 = -0.992$, $3.6\sigma$ DESI tension
    (Sim~109).  Full $m_0$ scan confirms structural obstruction:
    $\rho_\Lambda$ dominates at all Planck-consistent $m_0$ (Sim~111).
  \item Dark matter: all kinematic modifications exhausted
    (Sims~99--100, 103).  Coupling $\beta\sim 1600\,\Lz$ required but
    absent from locked action.
\end{itemize}

\subsection{The two structural gaps}

The locked action \eqref{eq:locked} contains two unexplained external
inputs:

\paragraph{Gap 1: the bare cosmological constant.}
The Friedmann equation includes $\rho_\Lambda = \Omega_\Lambda \times
3H_0^2$ by hand.  The $\Psi$ field contributes a correction of order
$\Lz \ll 1$ to the total energy budget and cannot generate $\rho_\Lambda$
dynamically.  Sim~111 confirmed that varying $m_0$ cannot change this: as
long as $\rho_\Lambda$ is present and dominant, $w\approx -1$ universally.

\paragraph{Gap 2: no matter coupling.}
$S_{\rm matter}$ appears in \eqref{eq:locked} without any coupling
to $\Psi$.  The $\Psi$--curvature coupling $\Lz\Psi^2 R$ provides an
indirect link (through the Einstein equation $R \sim T$), but it is
suppressed by $\Lz \ll 1$.  Sims~99--100 and 103 showed this is
insufficient to produce galactic rotation curves by factors of
$10^6$--$10^7$.

Both gaps require the action to be \emph{unlocked}.

\subsection{Strategy}

We unlock the action by deriving the \emph{unique} extensions permitted
by the symmetry and consistency principles that defined Phase~1 in the
first place.  No new principles are introduced; the Phase~2 action is
fully determined by applying existing principles more completely.

% ============================================================
\section{Symmetry Principles of RIFT}
\label{sec:principles}

The following principles were either explicit in Phase~1 or are
necessary consequences of it.

\subsection{Diffeomorphism invariance}
The action must be a scalar under general coordinate transformations.
This is preserved by the locked action and all extensions below.

\subsection{\texorpdfstring{$\mathbb{Z}_2$}{Z2} symmetry}
The field equation
\begin{equation}
  \Box\Psi + m_0^2\Psi = 2\Lz\Psi R
  \label{eq:kg}
\end{equation}
is invariant under $\Psi\to-\Psi$.  This symmetry must be preserved
in the Phase~2 action: all new terms must be even in $\Psi$.
Physically, the sign of the memory field is unobservable; only $\Psi^2$
couples to geometry.

\subsection{Horndeski class (second-order field equations)}
The Phase~1 action is in the Horndeski class.  Sim~108 confirmed
$c_T = c$ exactly, which is protected by the $G_5 = 0$ Horndeski
sub-class.  We require field equations to remain second-order to
avoid Ostrogradsky ghosts.

\subsection{UV finiteness}
Sims~102--106 established that the locked action is UV finite through
two loops under Dyson resummation.  The Phase~2 extensions must not
introduce new UV divergences beyond those already regulated.
Specifically: the matter coupling $\beta\Psi^2\mathcal{L}_m$ introduces
a new vertex $\beta\Psi^2 T_{\mu\nu}$; its one-loop contribution to
$\Psi$ self-energy is $\sim\beta^2 T^2/k_m^4$, regulated by the same
memory cutoff.

\subsection{The recursion principle}
This is the defining principle of RIFT: \emph{$\Psi$ encodes the
memory of curvature and is sourced by it}.  The Klein-Gordon equation
\eqref{eq:kg} expresses this.  The principle implies:

\begin{enumerate}
  \item[\textbf{R1.}] The source of $\Psi$ is $R$ (curvature), not $T$
    (matter) directly.  (This was locked in Phase~1.)
  \item[\textbf{R2.}] The coupling of $\Psi$ to matter is \emph{mediated}
    by curvature.  Matter enters through the Einstein equation
    $R\sim T + T_\Psi$.  Any direct $\Psi$--matter coupling must be
    consistent with this mediation.
  \item[\textbf{R3.}] The equilibrium of $\Psi$ is a \emph{fixed point}
    of the recursion: the field value that self-consistently sources the
    curvature that sustains it.
\end{enumerate}

Principle R3 is new in Phase~2 and is the key to deriving $\beta$.

% ============================================================
\section{Derivation of the Phase~2 Action}
\label{sec:action}

\subsection{Allowed terms}

Given the five principles above, the most general Phase~2 action is:

\begin{equation}
S = \int d^4x\,\sqrt{-g}\left[
  \frac{\pl^2 + 2\Lz\Psi^2}{2}\,R
  - \half(\nabla\Psi)^2
  - V_J(\Psi)
  + \beta\Psi^2\mathcal{L}_m^{(\rm DM)}
\right]
+ S_{\rm SM}
\label{eq:phase2}
\end{equation}

where:
\begin{itemize}
  \item $V_J(\Psi)$ is the Jordan-frame scalar potential, to be derived.
  \item $\beta\Psi^2\mathcal{L}_m^{(\rm DM)}$ is the conformal matter
    coupling.  The $\mathbb{Z}_2$ symmetry requires it to be even in
    $\Psi$; the lowest-order term consistent with Horndeski is
    $\beta\Psi^2\mathcal{L}_m$.
  \item $S_{\rm SM}$ is the Standard Model action, uncoupled to $\Psi$
    (baryons and photons remain on GR geodesics).
  \item \textbf{No bare cosmological constant.}  $\rho_\Lambda$ is
    absent; it must be reproduced by $V_J$.
\end{itemize}

\noindent The Phase~2 action has two new parameters: $\lambda$ (in $V_J$)
and $\beta$.  Both are determined by observational boundary conditions.

\subsection{Deriving \texorpdfstring{$V_J(\Psi)$}{VJ(Psi)} — dark energy}
\label{sec:vj}

\paragraph{Requirement.}
We require that the action reproduces the observed late-time
acceleration without a bare $\Lambda$.  That is, the effective
dark-energy density today must satisfy $\rho_{\rm DE}(z=0)\approx
0.685\times 3H_0^2\,\pl^2$.

\paragraph{Einstein-frame transformation.}
Under the conformal rescaling $\tilde{g}_{\mu\nu} = F(\Psi)\,g_{\mu\nu}$
with $F = (\pl^2 + 2\Lz\Psi^2)/\pl^2$, the Jordan-frame action maps to:

\begin{equation}
S_E = \int d^4x\,\sqrt{-\tilde{g}}\left[
  \frac{\pl^2}{2}\tilde{R}
  - \half K(\Psi)(\tilde\nabla\Psi)^2
  - V_E(\Psi)
\right]
+ \tilde{S}_{\rm matter}
\label{eq:einstein}
\end{equation}

where the Einstein-frame potential is:
\begin{equation}
\boxed{V_E(\Psi) = \frac{V_J(\Psi)}{F(\Psi)^2}
= \frac{V_J(\Psi)}{\left(1 + 2\Lz\Psi^2/\pl^2\right)^2}}
\label{eq:ve}
\end{equation}

and the kinetic factor $K(\Psi)$ modifies the canonical normalisation
of $\Psi$ in the Einstein frame.

\paragraph{Why a quartic potential?}
Consider the hierarchy of $\mathbb{Z}_2$-symmetric potentials:
\begin{align}
  V_J^{(2)} &= \half m_0^2 \Psi^2 \quad\text{[locked action]}\\
  V_J^{(4)} &= \lambda\Psi^4 \quad\text{[quartic extension]}
\end{align}

For $V_J^{(2)}$: $V_E(\Psi\to\infty)\to 0$.  The Einstein-frame
potential goes to zero at large field values.  This cannot sustain a
de Sitter phase.

For $V_J^{(4)} = \lambda\Psi^4$:
\begin{equation}
V_E(\Psi) = \frac{\lambda\Psi^4}{\left(1+2\Lz\Psi^2/\pl^2\right)^2}
\quad\xrightarrow{\Psi\gg\pl/\sqrt{2\Lz}}\quad
\frac{\lambda\pl^4}{4\Lz^2} \equiv \rho_{\rm plateau}
\label{eq:plateau}
\end{equation}

The quartic Jordan-frame potential generates a \textbf{plateau} in the
Einstein frame --- an effective cosmological constant
$\rho_{\rm plateau} = \lambda\pl^4/(4\Lz^2)$.  This is the
Starobinsky/Higgs-inflation mechanism, applied here to dark energy
rather than inflation.

The DE scale requirement fixes:
\begin{equation}
\lambda = \frac{4\Lz^2\,\rho_{\rm DE,0}}{\pl^4}
= 4\Lz^2\,\Omega_\Lambda\frac{3H_0^2}{\pl^2}
\approx 4\times(0.003)^2\times 0.685\times\frac{3H_0^2}{\pl^2}
\label{eq:lambda_fix}
\end{equation}

In natural units $H_0 = 1$, $8\pi G = 1$ (as in Phase~1 simulations):
$\pl^2 = 1/(8\pi G) = 1$ and $\rho_{\rm DE,0} = 0.685\times 3 = 2.055$.
This gives:
\begin{equation}
\lambda = 4\times(0.003)^2\times 2.055 \approx 7.4\times10^{-5}
\label{eq:lambda_value}
\end{equation}

\paragraph{The slow-roll condition for DESI.}
For the DESI $w_0$--$w_a$ hint to be reproduced, $\Psi$ must be
\emph{rolling today}, not frozen at the plateau.  The field rolls
when the slow-roll parameter
\begin{equation}
\epsilon_V = \frac{\pl^2}{2}\left(\frac{V_E'}{V_E}\right)^2
\end{equation}
is $O(1)$ at $z=0$.  The transition from slow-roll ($\epsilon_V\ll1$)
to rolling ($\epsilon_V\sim1$) occurs at the inflection point of $V_E$,
which lies at $\Psi_*\approx\pl/\sqrt{2\Lz}$.

If $\Psi_0\lesssim\Psi_*$ today, the field is on the \emph{slope} of
$V_E$, giving $w>-1$ (thawing quintessence).  This is the regime that
resolves the DESI tension.  SIM112 will determine the value of $\Psi_0$
that simultaneously satisfies the Friedmann constraint and the DESI
$w_0$--$w_a$ posterior.

\subsection{Deriving \texorpdfstring{$\beta$}{beta} — dark matter}
\label{sec:beta}

\paragraph{Effective mass in matter.}
The conformal coupling $\beta\Psi^2\mathcal{L}_m$ modifies the
Klein-Gordon equation. For pressureless matter ($\mathcal{L}_m = -\rho_m$
in the fluid limit):
\begin{equation}
\Box\Psi + m_{\rm eff}^2(\mathbf{x})\,\Psi = 2\Lz\Psi R
\label{eq:kg_eff}
\end{equation}
with
\begin{equation}
m_{\rm eff}^2(\mathbf{x}) = m_0^2 - 2\beta\rho_m(\mathbf{x})
\label{eq:meff}
\end{equation}

In vacuum ($\rho_m = 0$): $m_{\rm eff}^2 = m_0^2 > 0$ (stable).
In a dense region: if $2\beta\rho_m > m_0^2$, then $m_{\rm eff}^2 < 0$
--- a \textbf{tachyonic instability}.  The field grows until nonlinear
terms (from $V_J = \lambda\Psi^4$) stabilise it at a new condensate
value $\Psi_c$.

\paragraph{The condensate value.}
Setting $\Box\Psi = 0$ (static, uniform condensate) and $R\approx -8\pi G T\approx 8\pi G\rho_m$ (matter-dominated halo), the fixed-point equation is:
\begin{equation}
m_0^2\Psi_c - 2\beta\rho_m\Psi_c + 4\lambda\Psi_c^3 = 2\Lz\Psi_c R
\label{eq:fixed_point}
\end{equation}

Dropping the small $m_0^2\Psi_c$ term and using $R = 8\pi G\rho_m$:
\begin{equation}
4\lambda\Psi_c^2 = (2\beta - 16\pi G\Lz)\,\rho_m
\label{eq:condensate}
\end{equation}

The condensate energy density is:
\begin{equation}
\rho_\Psi^{\rm cond} \approx \lambda\Psi_c^4 = \frac{(2\beta - 16\pi G\Lz)^2\rho_m^2}{64\lambda}
\label{eq:rho_cond}
\end{equation}

\paragraph{The recursion fixed-point condition.}
Principle R3 requires that $\rho_\Psi^{\rm cond}$ self-consistently
sources the curvature that drives the condensation.  In a dark matter
halo, the total effective DM density is $\rho_{\rm DM} = \rho_\Psi^{\rm cond}$,
and the observed rotation velocity gives:
\begin{equation}
v_c^2(r) = \frac{G M_{\rm DM}(<r)}{r}
\end{equation}
For a flat rotation curve at $v_c \approx 200\,{\rm km/s}$ in a
Milky-Way-like galaxy, with typical halo density
$\rho_{\rm DM} \approx 0.3\,{\rm GeV/cm}^3$ at the Solar radius:

\begin{equation}
\beta = \half\left(\frac{8\sqrt{\lambda}\,\rho_{\rm DM}}{\rho_m^{}} + 16\pi G\Lz\right)
\label{eq:beta_fix}
\end{equation}

Using $\lambda\approx 7.4\times10^{-5}$, $\rho_{\rm DM}/\rho_m\approx 5$
(local DM-to-baryon ratio), and $16\pi G\Lz\ll\beta$:
\begin{equation}
\beta \approx 4\sqrt{\lambda} \approx 4\times(7.4\times10^{-5})^{1/2} \approx 0.034
\label{eq:beta_estimate}
\end{equation}

Note: $\beta\approx 0.034\approx 11\,\Lz$ --- within an order of
magnitude of $\Lz$, not the factor-of-1600 discrepancy found in the
locked action.  The quartic potential changes the algebra entirely by
providing a new stabilisation scale.  SIM113 will test whether this
estimate produces correct rotation curve profiles.

\subsection{Field equations of the Phase~2 action}

Varying \eqref{eq:phase2} with respect to $g^{\mu\nu}$ and $\Psi$:

\paragraph{Modified Einstein equation.}
\begin{multline}
(\pl^2 + 2\Lz\Psi^2)G_{\mu\nu}
= T_{\mu\nu}^{(\rm SM)} + T_{\mu\nu}^{(\rm DM)}
+ \partial_\mu\Psi\partial_\nu\Psi
- g_{\mu\nu}\left[\half(\nabla\Psi)^2 + V_J\right]\\
+ 2\Lz\left[g_{\mu\nu}\Box(\Psi^2) - \nabla_\mu\nabla_\nu(\Psi^2)\right]
\label{eq:einstein_eq}
\end{multline}

\paragraph{Modified Klein-Gordon equation.}
\begin{equation}
\Box\Psi + 4\lambda\Psi^3 - 2\beta\rho_m\Psi = 2\Lz\Psi R
\label{eq:kg_phase2}
\end{equation}

The $m_0^2\Psi$ term from the locked action is dropped; it is replaced
by $4\lambda\Psi^3$ (the quartic stabiliser) and $-2\beta\rho_m\Psi$
(the condensation driver).

Note: the locked-action Klein-Gordon equation \eqref{eq:kg} is recovered
in the limit $\lambda\to 0$, $\beta\to 0$, $m_0^2>0$.  Phase~2 is a
strict superset of Phase~1.

% ============================================================
\section{Physical Interpretation}
\label{sec:physics}

\subsection{Dark energy as geometric quintessence}

The key conceptual shift is that dark energy is no longer an
unexplained constant added to the stress-energy tensor.  It is the
energy stored in the $\Psi$ field as it rolls down the Einstein-frame
plateau potential \eqref{eq:ve}.

At early times ($z\gg 1$, matter domination): $R$ is large; Ψ is
driven to large values by the source term $2\Lz\Psi R$; $V_E$ is near
the plateau; the field is nearly frozen ($w\approx -1$).

At late times ($z\lesssim 1$, dark energy domination): $R\to -4H^2$
(de Sitter); the source weakens; $\Psi$ begins to roll; $w$ deviates
from $-1$ toward the DESI regime ($w>-1$).

This is \emph{thawing quintessence} with a \emph{dynamically
generated plateau} --- not inserted, but derived from the
$\Lz\Psi^2 R$ non-minimal coupling.

\subsection{Dark matter as curvature condensate}

In galaxies, the combination of matter coupling $\beta\rho_m$ and
curvature coupling $2\Lz R$ drives $m_{\rm eff}^2 < 0$ in dense
regions.  The field condenses to a non-zero value $\Psi_c$ stabilised
by the quartic $\lambda\Psi^4$ term.

The condensate $\rho_\Psi^{\rm cond}\propto\rho_m^2$ has a
density profile that traces the baryon density squared --- a
non-trivial prediction distinguishable from NFW or isothermal profiles.
At large $r$ (where $\rho_m\to 0$), the condensate decays as
$\rho_\Psi\propto\rho_m^2\to 0$.  The condensate is \emph{baryon-seeded}:
there is no dark matter without baryons.

This is testable: RIFT predicts a $\rho_{\rm DM}\propto\rho_m^2$
profile at galactic scales, distinct from the $\Lambda$CDM NFW profile.

\subsection{Single-parameter connection between DE and DM}

Note that both $\lambda$ (DE scale, eq.~\ref{eq:lambda_fix}) and
$\beta$ (DM coupling, eq.~\ref{eq:beta_estimate}) are determined by
$\Lz$ and observational inputs.  Remarkably:
\begin{equation}
\beta \approx 4\sqrt{\lambda} = 4\sqrt{\frac{4\Lz^2\,\rho_{\rm DE,0}}{\pl^4}}
= \frac{8\Lz\sqrt{\rho_{\rm DE,0}}}{\pl^2}
\label{eq:connection}
\end{equation}

If this relation holds exactly, then $\beta$ and $\lambda$ are
\emph{not independent}: both are functions of $\Lz$ alone.  The
Phase~1 parameter $\Lz = 0.003$ would then determine the entire
Phase~2 phenomenology.  SIM112--114 will test whether the
fixed-point condition \eqref{eq:beta_fix} is consistent with
$\beta\approx 8\Lz\sqrt{\rho_{\rm DE,0}}/\pl^2$.

% ============================================================
\section{Phase~2 Simulation Programme}
\label{sec:sims}

The following simulations will test the Phase~2 action.

\begin{table}[h]
\centering
\caption{Phase~2 simulation programme.}
\begin{tabular}{cll}
\toprule
Sim & Description & Pass criterion \\
\midrule
SIM112 & Remove $\rho_\Lambda$; add $\lambda\Psi^4$.
  Solve background. Find $\lambda$ s.t.\
  $\Omega_{\rm DE}(z=0)=0.685$.
  Refit $w_0$--$w_a$.
  & $w_0$ within $2\sigma$ Planck; \\
  &&  DESI tension $<2\sigma$ \\
\midrule
SIM113 & Add $\beta\Psi^2\rho_m$ coupling.
  Solve galactic Ψ profile (static + time-domain).
  Test rotation curves.
  & Flat $v_c(r)$ reproduced \\
  & & at $r\in[2,30]$~kpc \\
\midrule
SIM114 & Joint cosmological test.
  Run Phase~1 SIM88--98 suite
  with Phase~2 action.
  & $\Delta\chi^2 \leq 0$ vs Phase~1 \\
  & & in all 22 tests \\
\midrule
SIM115 & Solar System screening with $\beta$ coupling.
  Check fifth-force bounds.
  & $|\gamma-1| < 2.3\times10^{-5}$ \\
\bottomrule
\end{tabular}
\label{tab:sims}
\end{table}

\noindent The simulations will be run in order.  If SIM112 fails (cannot
reproduce $\Omega_{\rm DE}$ without fine-tuning $\lambda$ beyond the
estimate \eqref{eq:lambda_value}), the action is revised before
proceeding.  If SIM113 fails (rotation curves wrong), $\beta$ is
adjusted and the fixed-point condition is re-examined.

% ============================================================
\section{Open Theoretical Questions}
\label{sec:open}

\subsection{Mass term}

The locked action had $\half m_0^2\Psi^2$.  In Phase~2 with $\lambda\Psi^4$,
the field's mass at the minimum of $V_J$ vanishes at $\Psi=0$.  A mass
term is generated by quantum corrections (SIM102: $\Sigma(0)=\Lz^2 k_m^4/(64\pi^2)$).
Whether this is large enough to matter for the DE or DM phenomenology
is a question for the simulations.

\subsection{Inflation compatibility}

SIM110 showed that $\xi\sim10^4$ is needed for the Starobinsky attractor.
The Phase~2 action, with large $\Psi$ during inflation and $\lambda\Psi^4$
as the potential, enters the Starobinsky regime at $\Psi\gg\pl/\sqrt{2\Lz}$.
The inflationary plateau from eq.~\eqref{eq:plateau} is the same mechanism
as the DE plateau, but at a much higher energy scale.  Whether a single
$\lambda$ can serve both inflation and DE simultaneously is an open
question; it likely requires a running $\lambda(\mu)$ from the RG flow.

\subsection{Baryon-seeded DM and structure formation}

The RIFT DM condensate $\rho_{\rm DM}\propto\rho_m^2$ predicts a very
different matter power spectrum from \lcdm{}.  At large scales (where
$\rho_m$ is smooth), RIFT DM is suppressed; at galactic scales, it is
concentrated.  This has potentially large consequences for the CMB
and BAO tests established in Phase~1 and must be checked in SIM114.

\subsection{The \texorpdfstring{$\beta$}{beta}--$\Lz$ relation}

Equation~\eqref{eq:connection} suggests $\beta = 8\Lz\sqrt{\rho_{\rm DE,0}}/\pl^2$
is exact.  If this relation is a consequence of the fixed-point condition
R3 (not an approximation), it would reduce the Phase~2 action to a
\emph{one-parameter theory} ($\Lz$ only) --- a remarkable unification
of DE and DM under a single coupling.  Deriving this relation
analytically is a priority theoretical task.

% ============================================================
\section{Conclusion}
\label{sec:conclusion}

We have derived the Phase~2 RIFT action from first principles.
Starting from the symmetry and consistency principles that defined
Phase~1 --- diffeomorphism invariance, $\mathbb{Z}_2$, Horndeski,
UV finiteness, and the recursion principle --- the unique extension
of the locked action is:

\begin{equation}
\boxed{
S = \int d^4x\,\sqrt{-g}\left[
  \frac{\pl^2 + 2\Lz\Psi^2}{2}\,R
  - \half(\nabla\Psi)^2
  - \lambda\Psi^4
  + \beta\Psi^2\mathcal{L}_m
\right]
+ S_{\rm SM}
}
\label{eq:phase2_final}
\end{equation}

with $\Lz = 0.003$ (locked), $\lambda\approx 7.4\times10^{-5}$
(DE scale), $\beta\approx 0.034$ (DM coupling), and \textbf{no bare
cosmological constant}.

The three key predictions to be tested in Phase~2 simulations are:
\begin{enumerate}
\item The quartic potential generates DE via an Einstein-frame plateau,
  giving thawing quintessence with $w_0 > -1$ (DESI regime).
\item The conformal matter coupling drives Ψ condensation in galaxies,
  generating a $\rho_{\rm DM}\propto\rho_m^2$ dark matter profile.
\item All Phase~1 cosmological tests ($\Delta\chi^2\leq 0$) are
  preserved, since the Phase~2 action reduces to the locked action
  at $\lambda\to 0$, $\beta\to 0$.
\end{enumerate}

If these tests pass, RIFT Phase~2 will be the first theory in this
programme to explain dark energy and dark matter from a single
geometric principle, with no free parameters beyond those fixed by
observations.

% ============================================================
\end{document}
